\section{The Appeal of Tradition}

The appeal of Julius Evola to a segment of disaffected youth is undeniable. Disgusted with the superficial, egalitarian, and epicene culture of the West, and unable to find a suitable home in the existent spiritual organisations, Evola provides an alternative that is absolutely unique.

Once the miasma hovering over Western society is blown away and there is an opening through which one can see a little more clearly, the emotional attachment to the shallow type of person so admired today — politicians, movie and rock stars, professional athletes, emaciated models, industrial magnates, self-help gurus, and all the so-called experts — suddenly loses all its lustre and appeal.

Evola, too, despises this world and offers quite an alternative — the world of Tradition. This world is populated with kings, heroes, warriors, knights of chivalry, sages and wise men, magicians, alchemists, yogis, tantrists, clairvoyants, mystics with occult powers.

He tours the mythology of Traditional societies: from Samurais in Japan, to Vedic India, Mithraism of Iran, the gods and goddesses of pagan Greece and Rome, to the Valhalla of the Germanic and Norse tribes, even to the Aztecs of Mexico. From the commonalties, he reconstructs Traditional societies as they must have existed in prehistory: their hierarchical and caste-based social structures, their admiration for and cultivation of the ``higher" type of man.

For Evola, the Traditional world ended with the waning of the Catholic middle ages. Successive waves of revolution — the Renaissance, the Reformation, the Enlightenment, industrialism, democracy, socialism, feminism — in short, all the movements of thought that we today praise as the highest accomplishments of the West — represent, for him, mere decay and decadence.

He rejects science as providing only contingent truths and suitable only for the mass man of our time. He opposes it with Wisdom resulting from the attainment of higher consciousness, which represents an achievement of the individual.

Instead of constant class warfare, Evola envisions a stable hierarchy in which each man freely and willingly takes his place.

Evola is the rare type of man who follows his line of thought intransigently all the way through, wherever it leads him. He shows us the aetiology and symptoms of the spiritual malaise of the West. The cure he proposes is a different story and is essentially left unfinished. He calls for a new counter-Reformation that will go back through the Middle Ages to recover once again the values of Tradition. The man to do that is yet to appear.



\flrightit{Posted on 2007-11-27 by Cologero }
